原树邻接表类型为 \texttt{vector<vector<Edge>>}，其中
\texttt{Edge=pair<ll,int>}，按 \texttt{\{边权,终点\}} 存边。
构造 \texttt{VirtualTree T(g,root)} 后，调用 \texttt{T.build(key)} 建立包含所有关键点的最小虚树，并返回虚树根。

模板只排序一次关键点，再在线插入 LCA、弹栈并连边。
\texttt{T.vt[u]} 的元素也是 \texttt{\{权值,v\}}，虚树边权为原树路径权值和
\[
  dis[v]-dis[u].
\]
\texttt{dist(u,v)} 返回原树两点间路径权值和。

\begin{icpcresult}[多次构建]
  \texttt{nodes}、\texttt{stk} 和 \texttt{buf} 会跨询问复用容量；每次只清空上一次虚树涉及的邻接表。
  关键点允许重复，模板会自动去重，虚树节点数量不超过 $2k-1$。
\end{icpcresult}

\begin{icpcwarning}[nodes 顺序]
  为省去恢复 DFS 序的额外遍历，\texttt{nodes} 只保证包含本次虚树的全部节点，不保证 DFS 序或拓扑序。
  需要树形 DP 时应从 \texttt{build()} 返回的根沿 \texttt{vt} 遍历。
\end{icpcwarning}

\begin{icpcwarning}[根与数值范围]
  模板不会自动加入原树根；若 DP 必须从原树根开始，应把它加入关键点。
  根路径权值和与虚树边权使用 \texttt{ll}，仍需保证加减法不溢出。
\end{icpcwarning}
